\par\SourceLeaf{428}%
\noindent\hspace*{1.2em}Da man durch Elimination aus der Gleichung $f(x,y,z)=0$ die eine Koordinate $z$ als Funktion von $x$ und $y$ darstellen kann, nehmen wir der Einfachheit halber an, dies sei geschehen und
\[
z=F(x,y)
\]
sei gefunden worden. Man kann dann als Gleichung der Fläche unmittelbar
\[
z-F(x,y)=0
\]
setzen oder
\[
f(x,y,z)=z-F(x,y).
\]

Setzen wir also
\[
dF(x,y)=tdx+udy
\]
so erhalten wir
\[
P=-t,\qquad Q=-u,\qquad R=1,
\]
wobei $t,u$ nur Funktionen von $x$ und $y$ sind. Wir setzen ferner
\[
dt=Tdx+Udy,\qquad du=Udx+Vdy.
\]

Auf der ganzen Fläche gilt somit
\[
dz=tdx+udy
\]
und daher auf der Kurve
\[
\zeta=t\xi+u\eta.
\]
Durch Differenzieren und Einsetzen der obigen Werte für $d\xi,d\eta,d\zeta$ erhält man
\[
\begin{aligned}
(\zeta'-t\xi'-u\eta')\frac{ds}{r}&=\xi dt+\eta du\\
&=(\xi\xi T+2\xi\eta U+\eta\eta V)ds
\end{aligned}
\]
oder
\[
\begin{aligned}
\frac1r&=\frac{\xi\xi T+2\xi\eta U+\eta\eta V}{-\xi't-\eta'u+\zeta'}\\
&=\frac{Z(\xi\xi T+2\xi\eta U+\eta\eta V)}{X\xi'+Y\eta'+Z\zeta'}\\
&=\frac{Z(\xi\xi T+2\xi\eta U+\eta\eta V)}{\cos L\lambda'}.
\end{aligned}
\]
